% Options for packages loaded elsewhere
% Options for packages loaded elsewhere
\PassOptionsToPackage{unicode}{hyperref}
\PassOptionsToPackage{hyphens}{url}
\PassOptionsToPackage{dvipsnames,svgnames,x11names}{xcolor}
%
\documentclass[
  letterpaper,
  DIV=11,
  numbers=noendperiod]{scrreprt}
\usepackage{xcolor}
\usepackage{amsmath,amssymb}
\setcounter{secnumdepth}{5}
\usepackage{iftex}
\ifPDFTeX
  \usepackage[T1]{fontenc}
  \usepackage[utf8]{inputenc}
  \usepackage{textcomp} % provide euro and other symbols
\else % if luatex or xetex
  \usepackage{unicode-math} % this also loads fontspec
  \defaultfontfeatures{Scale=MatchLowercase}
  \defaultfontfeatures[\rmfamily]{Ligatures=TeX,Scale=1}
\fi
\usepackage{lmodern}
\ifPDFTeX\else
  % xetex/luatex font selection
\fi
% Use upquote if available, for straight quotes in verbatim environments
\IfFileExists{upquote.sty}{\usepackage{upquote}}{}
\IfFileExists{microtype.sty}{% use microtype if available
  \usepackage[]{microtype}
  \UseMicrotypeSet[protrusion]{basicmath} % disable protrusion for tt fonts
}{}
\makeatletter
\@ifundefined{KOMAClassName}{% if non-KOMA class
  \IfFileExists{parskip.sty}{%
    \usepackage{parskip}
  }{% else
    \setlength{\parindent}{0pt}
    \setlength{\parskip}{6pt plus 2pt minus 1pt}}
}{% if KOMA class
  \KOMAoptions{parskip=half}}
\makeatother
% Make \paragraph and \subparagraph free-standing
\makeatletter
\ifx\paragraph\undefined\else
  \let\oldparagraph\paragraph
  \renewcommand{\paragraph}{
    \@ifstar
      \xxxParagraphStar
      \xxxParagraphNoStar
  }
  \newcommand{\xxxParagraphStar}[1]{\oldparagraph*{#1}\mbox{}}
  \newcommand{\xxxParagraphNoStar}[1]{\oldparagraph{#1}\mbox{}}
\fi
\ifx\subparagraph\undefined\else
  \let\oldsubparagraph\subparagraph
  \renewcommand{\subparagraph}{
    \@ifstar
      \xxxSubParagraphStar
      \xxxSubParagraphNoStar
  }
  \newcommand{\xxxSubParagraphStar}[1]{\oldsubparagraph*{#1}\mbox{}}
  \newcommand{\xxxSubParagraphNoStar}[1]{\oldsubparagraph{#1}\mbox{}}
\fi
\makeatother


\usepackage{longtable,booktabs,array}
\usepackage{calc} % for calculating minipage widths
% Correct order of tables after \paragraph or \subparagraph
\usepackage{etoolbox}
\makeatletter
\patchcmd\longtable{\par}{\if@noskipsec\mbox{}\fi\par}{}{}
\makeatother
% Allow footnotes in longtable head/foot
\IfFileExists{footnotehyper.sty}{\usepackage{footnotehyper}}{\usepackage{footnote}}
\makesavenoteenv{longtable}
\usepackage{graphicx}
\makeatletter
\newsavebox\pandoc@box
\newcommand*\pandocbounded[1]{% scales image to fit in text height/width
  \sbox\pandoc@box{#1}%
  \Gscale@div\@tempa{\textheight}{\dimexpr\ht\pandoc@box+\dp\pandoc@box\relax}%
  \Gscale@div\@tempb{\linewidth}{\wd\pandoc@box}%
  \ifdim\@tempb\p@<\@tempa\p@\let\@tempa\@tempb\fi% select the smaller of both
  \ifdim\@tempa\p@<\p@\scalebox{\@tempa}{\usebox\pandoc@box}%
  \else\usebox{\pandoc@box}%
  \fi%
}
% Set default figure placement to htbp
\def\fps@figure{htbp}
\makeatother

\ifLuaTeX
  \usepackage{luacolor}
  \usepackage[soul]{lua-ul}
\else
  \usepackage{soul}
\fi




\setlength{\emergencystretch}{3em} % prevent overfull lines

\providecommand{\tightlist}{%
  \setlength{\itemsep}{0pt}\setlength{\parskip}{0pt}}



 


\KOMAoption{captions}{tableheading}
\makeatletter
\@ifpackageloaded{bookmark}{}{\usepackage{bookmark}}
\makeatother
\makeatletter
\@ifpackageloaded{caption}{}{\usepackage{caption}}
\AtBeginDocument{%
\ifdefined\contentsname
  \renewcommand*\contentsname{Table of contents}
\else
  \newcommand\contentsname{Table of contents}
\fi
\ifdefined\listfigurename
  \renewcommand*\listfigurename{List of Figures}
\else
  \newcommand\listfigurename{List of Figures}
\fi
\ifdefined\listtablename
  \renewcommand*\listtablename{List of Tables}
\else
  \newcommand\listtablename{List of Tables}
\fi
\ifdefined\figurename
  \renewcommand*\figurename{Figure}
\else
  \newcommand\figurename{Figure}
\fi
\ifdefined\tablename
  \renewcommand*\tablename{Table}
\else
  \newcommand\tablename{Table}
\fi
}
\@ifpackageloaded{float}{}{\usepackage{float}}
\floatstyle{ruled}
\@ifundefined{c@chapter}{\newfloat{codelisting}{h}{lop}}{\newfloat{codelisting}{h}{lop}[chapter]}
\floatname{codelisting}{Listing}
\newcommand*\listoflistings{\listof{codelisting}{List of Listings}}
\makeatother
\makeatletter
\makeatother
\makeatletter
\@ifpackageloaded{caption}{}{\usepackage{caption}}
\@ifpackageloaded{subcaption}{}{\usepackage{subcaption}}
\makeatother
\usepackage{bookmark}
\IfFileExists{xurl.sty}{\usepackage{xurl}}{} % add URL line breaks if available
\urlstyle{same}
\hypersetup{
  pdftitle={Guia\_Rapida\_MarckDown},
  pdfauthor={Tec. Sup. Christopher Jara},
  colorlinks=true,
  linkcolor={blue},
  filecolor={Maroon},
  citecolor={Blue},
  urlcolor={Blue},
  pdfcreator={LaTeX via pandoc}}


\title{Guia\_Rapida\_MarckDown}
\author{Tec. Sup. Christopher Jara}
\date{2027-04-10}
\begin{document}
\maketitle

\renewcommand*\contentsname{Table of contents}
{
\hypersetup{linkcolor=}
\setcounter{tocdepth}{2}
\tableofcontents
}

\bookmarksetup{startatroot}

\chapter*{Prefacio}\label{prefacio}
\addcontentsline{toc}{chapter}{Prefacio}

\markboth{Prefacio}{Prefacio}

¡Bienvenidos al \textbf{Manual Práctico de MarkDown}!

Este libro digital ha sido preparado como material de estudio y guía de
referencia para ustedes, los estudiantes de tercero de bachillerato de
la Unidad Educativa ``Carlos Crespi II''.

\section*{¿Por qué aprender
Markdown?}\label{por-quuxe9-aprender-markdown}
\addcontentsline{toc}{section}{¿Por qué aprender Markdown?}

\markright{¿Por qué aprender Markdown?}

En el mundo digital actual, la forma en que escribimos y compartimos
información está en constante evolución. Markdown es un lenguaje que nos
permite dar formato a nuestros textos (como títulos, negritas, listas e
imágenes) de una manera increíblemente simple y eficiente, usando solo
texto plano.

Dominar Markdown no solo les facilitará la toma de apuntes y la
redacción de documentos técnicos, sino que también es una habilidad
fundamental en el desarrollo de software (para crear la documentación de
proyectos), la publicación de blogs y la colaboración en plataformas
como GitHub. Es una herramienta que les servirá tanto en sus estudios
actuales como en su futura carrera profesional y universitaria.

\section*{¿Para quién es este libro?}\label{para-quiuxe9n-es-este-libro}
\addcontentsline{toc}{section}{¿Para quién es este libro?}

\markright{¿Para quién es este libro?}

Este manual está diseñado pensando en ustedes. No se requiere
experiencia previa. Empezaremos desde los conceptos más básicos,
asegurando que tengan una base sólida, antes de avanzar hacia temas más
complejos.

\section*{Estructura de este Libro}\label{estructura-de-este-libro}
\addcontentsline{toc}{section}{Estructura de este Libro}

\markright{Estructura de este Libro}

Este es un proyecto en desarrollo. El objetivo es construir, junto con
ustedes, un recurso completo. Por ahora, el libro se estructura de la
siguiente manera:

\begin{itemize}
\item
  \textbf{Parte I: Fundamentos de Markdown.} Comenzamos con un análisis
  detallado de la sintaxis esencial, cubriendo el estándar
  ``CommonMark'' y las útiles extensiones de ``GitHub Flavored
  Markdown'' (GFM). Esta primera parte les dará todas las herramientas
  que necesitan para el 90\% de sus tareas diarias.
\item
  \textbf{Partes Futuras (Próximamente).} A medida que avancemos,
  añadiremos nuevos capítulos. Estos explorarán temas avanzados como:
\item
  Creación de diagramas y flujos de trabajo.
\item
  Escritura de fórmulas matemáticas (LaTeX).
\item
  Aplicaciones prácticas de Markdown en sus proyectos académicos.
\item
  Introducción a herramientas de publicación como Quarto.
\end{itemize}

Espero que encuentren este manual útil y práctico. ¡Comencemos!

\bookmarksetup{startatroot}

\chapter{Introducción al Markdown}\label{introducciuxf3n-al-markdown}

Markdown se define como un lenguaje de marcado ligero diseñado para ser
fácil de leer y escribir en texto plano, con la capacidad de convertirse
sencillamente a HTML.

Sintaxis: Utiliza caracteres simples y fáciles de recordar (como \#, *,
{[}{]}).

Archivos: Los documentos son archivos de texto plano, usualmente con la
extensión .md.

Usos Comunes: Se emplea extensamente en GitHub, documentación técnica,
blogs, foros y para interactuar con sistemas de IA.

\section{Sintaxis Básica
(CommonMark)}\label{sintaxis-buxe1sica-commonmark}

Esta sección cubre los elementos fundamentales del Markdown original,
que hoy en día se ha estandarizado en gran medida como ``CommonMark''.

Encabezados: Se admiten seis niveles jerárquicos, definidos por el
número de almohadillas (\#) al inicio de la línea (ej.

es \# Título 1,

es \#\# Título 2).

Párrafos y Saltos: Los párrafos se separan dejando una línea en blanco
entre ellos. Para forzar un salto de línea simple (un ), se deben dejar
dos espacios en blanco al final de una línea antes de presionar enter.

Énfasis:

Negrita: Se logra envolviendo el texto con dobles asteriscos
(\textbf{texto en negrita}).

Cursiva: Se logra envolviendo el texto con asteriscos simples
(\emph{texto en cursiva}).

Reglas Horizontales: Se crean usando tres o más asteriscos (***) en una
línea separada.

Listas:

No Ordenadas: Se inician con -, * o +.

Ordenadas: Se inician con un número seguido de un punto (ej. 1.
Elemento).

Ambos tipos permiten anidación (sub-listas) mediante indentación.

Enlaces (Hipervínculos): La sintaxis principal es
\href{https://url.com}{Texto del enlace}.

Imágenes: Similar a los enlaces, pero con una exclamación al inicio:
\pandocbounded{\includegraphics[keepaspectratio]{url/imagen.jpg}}.

Citas (Blockquotes): Se utiliza el símbolo \textgreater{} al inicio de
la línea. Se pueden anidar usando múltiples símbolos (ej.
\textgreater\textgreater).

Código:

En línea (inline): Se envuelve el código con comillas invertidas simples
(\texttt{código}).

Bloques de código: Se utilizan comillas invertidas triples (```) para
delimitar bloques multilínea .

Escapado de Caracteres: Se usa la barra invertida () antes de un
carácter especial de Markdown (ej. *) para evitar que se interprete como
formato.

HTML en Markdown: Se permite incluir código HTML directamente dentro del
documento, lo cual es útil para elementos complejos.

\section{Sintaxis Extendida de GitHub
(GFM)}\label{sintaxis-extendida-de-github-gfm}

GitHub Flavored Markdown (GFM) es una extensión de CommonMark que añade
funcionalidades útiles, especialmente para desarrolladores.

Tablas: Se crean usando barras verticales (\textbar) para separar
columnas y guiones (-) para separar la fila del encabezado. Se puede
controlar la alineación del texto usando dos puntos (:) en la fila de
guiones (ej. :---: para centrar).

Tachado: Se envuelve el texto con dobles virgulillas (\st{texto
tachado}).

Listas de Tareas: Permiten crear casillas de verificación. Se usa {[}
{]} para una tarea pendiente y {[}x{]} para una completada.

Referencias y Menciones (GitHub): Atajos específicos de GitHub, como
(\textbf{menciones?}) de usuarios y \#número para referenciar Issues o
Pull Requests.

Notas a pie de página: Permiten añadir aclaraciones. Se usa {[}\^{}1{]}
en el texto y luego se define la nota con {[}\^{}1{]}: Texto de la
nota..

Alertas: Una extensión de las citas (blockquotes) para destacar
información. Usan una sintaxis especial como \textgreater{} {[}!NOTE{]}
o \textgreater{} {[}!WARNING{]} para crear bloques visuales distintivos.

\section{Contenido Avanzado}\label{contenido-avanzado}

La guía también menciona sintaxis avanzadas que muchas plataformas,
incluido GitHub (y crucialmente, Quarto), han adoptado.

LaTeX: Se utiliza para renderizar ecuaciones matemáticas complejas. Se
pueden usar \[...\] para bloques de ecuaciones.

Mermaid: Permite crear diagramas y flujos de trabajo (como diagramas de
flujo) directamente en Markdown. Esto se logra escribiendo la sintaxis
de Mermaid dentro de un bloque de código cercado e identificando el
lenguaje como mermaid.

\section{Herramientas y Recursos
Adicionales}\label{herramientas-y-recursos-adicionales}

Hoja de Trucos (Cheatsheet): La página 19 ofrece una tabla resumen muy
útil que contrasta la sintaxis clásica y la de GFM.

Herramientas Recomendadas: La guía concluye sugiriendo editores de texto
y herramientas. Es notable que Visual Studio Code es la primera
recomendación, lo cual refuerza su idoneidad para nuestro flujo de
trabajo con Quarto. También se mencionan Obsidian, Ghostwriter y iA
Writer.

Conclusión (Perspectiva de Ingeniero Quarto) Este documento es una
referencia fundamental. Para un desarrollador que trabaja con Quarto en
VS Code:

La Base es Sólida: Toda la ``Sintaxis Básica'' y la ``Sintaxis GFM''
(tablas, listas de tareas, notas a pie) son usadas directamente en
Quarto.

Alertas vs.~Callouts: Las ``Alertas'' de GFM son el precursor directo de
los ``Callouts'' de Quarto (ej. \{.callout-note\} \ldots{} ), que usamos
constantemente para estructurar documentos técnicos.

Contenido Avanzado Nativo: LaTeX y Mermaid son capacidades nativas de
Quarto. Saber que se pueden integrar es clave.

Entorno de Trabajo: La recomendación de VS Code alinea esta guía
perfectamente con nuestro entorno de desarrollo preferido.

\bookmarksetup{startatroot}

\chapter{Summary}\label{summary}

In summary, this book has no content whatsoever.

\bookmarksetup{startatroot}

\chapter*{References}\label{references}
\addcontentsline{toc}{chapter}{References}

\markboth{References}{References}

\phantomsection\label{refs}




\end{document}
